% ============================================================================
% Fichier  : annexes/Ann_D_textes.tex
% Titre    : Textes et Normes de Référence
% Version  : 1.0 -- novembre 2026
% ============================================================================
% Extraits identiques à ceux des chapitres. Loi n° 2010/012 : texte français
% officiel lu (Ann. D d'ACSI). COBAC R-2016/04 : art. 1 et 2 relevés sur une
% base en ligne, à collationner. Normes ISO : paraphrasées (texte sous
% licence) ; éditions vérifiées sur iso.org le 06/11/2026 pour 27000, 27001
% (et Amd 1:2024) et 27002.
% ============================================================================

\chapter{Textes et Normes de Référence}
\label{ann:textes}

Cette annexe rassemble les textes cités dans le cours, avec leur degré de
vérification. Seuls les extraits lus dans leur version officielle sont
reproduits entre guillemets (Chapitre~\ref{chap:droit}).

\section*{Degré de vérification}

\begin{tableaucours}{L{4.4cm}L{4.0cm}L{3.8cm}}{%
    \tableauHeader{Texte} &
    \tableauHeader{Degré de vérification} &
    \tableauHeader{Chapitres}
}
Loi n°~2010/012 (cybersécurité) & Texte officiel lu~; extraits cités &
\ref{chap:grc}, \ref{chap:iso27005}, \ref{chap:audit}, \ref{chap:droit},
\ref{chap:audit-technique} \\
Règlement COBAC R-2016/04 & Art.~1 et~2 relevés, à collationner &
\ref{chap:grc}, \ref{chap:droit} \\
Convention de Malabo (2014) & Présentée~; statut camerounais à vérifier &
\ref{chap:droit} \\
Loi n°~2024/017 (données personnelles) & À vérifier au Journal officiel &
\ref{chap:grc}, \ref{chap:droit}, \ref{chap:emergents} \\
Droit OHADA des sociétés & Présenté, à vérifier & \ref{chap:droit} \\
\end{tableaucours}

\section*{Loi n° 2010/012 relative à la cybersécurité et à la cybercriminalité}

\begin{boxjuris}
\textbf{Article~4, point~9}~: \textit{« Audit de sécurité~: examen
méthodique des composantes et des acteurs de la sécurité, de la politique,
des mesures, des solutions, des procédures et des moyens mis en oeuvre par
une organisation, pour sécuriser son environnement, effectuer des contrôles
de conformité, des contrôles d'évaluation de l'adéquation des moyens
(organisationnels, techniques, humains, financiers) investis au regard des
risques encourus, d'optimisation, de rationalité et de performance. »}

\medskip
\textbf{Article~7, alinéa~2} (extrait)~: l'Agence a notamment pour mission
\textit{« de s'assurer de la régularité, de l'effectivité des audits de
sécurité des systèmes d'information suivant les normes en la matière, des
organismes publics et des autorités de certification »}.

\medskip
\textbf{Article~13, alinéa~1}~: \textit{« Sont soumis à un audit de sécurité
obligatoire, les réseaux de communications électroniques et les systèmes
d'information, des opérateurs, des autorités de certification et des
fournisseurs de services de communications électroniques. »}

\medskip
\textbf{Article~14}~: \textit{« Le personnel de l'Agence et les experts
commis en vue d'accomplir des opérations d'audit sont astreints au secret
professionnel. »}

\medskip
\textbf{Article~26, alinéa~1}~: \textit{« Les exploitants des systèmes
d'information prennent toutes les mesures techniques et administratives afin
de garantir la sécurité des services offerts. A cet effet, ils se dotent de
systèmes normalisés leur permettant d'identifier, d'évaluer, de traiter et
de gérer continûment les risques liés à la sécurité des systèmes
d'information dans le cadre des services offerts directement ou
indirectement. »}

\medskip
\textbf{Article~29, alinéa~1}~: \textit{« Les exploitants des systèmes
d'information ont l'obligation de conserver les données de connexion et de
trafic de leurs systèmes d'information pendant une période de dix (10)
ans. »}

\medskip
\textbf{Article~32, alinéas~1 à~3}~: \textit{« (1) Les réseaux de
communications électroniques et les systèmes d'information sont soumis à un
régime d'audit de sécurité obligatoire et périodique de leurs systèmes de
sécurité par l'Agence. (2) L'audit de sécurité et les mesures d'impact de
gravité sont effectués chaque année ou lorsque les circonstances
l'exigent. (3) Les rapports d'audit sont confidentiels et adressés au
Ministre chargé des Télécommunications. »}

\medskip
\textbf{Article~61, alinéa~1}~: \textit{« Sont punis d'un emprisonnement de
trois (03) mois à trois (03) ans et d'une amende de 20.000 (vingt mille) à
100.000 (cent mille) F CFA, les personnels de l'Agence et les experts des
personnes morales chargés des audits qui révèlent sans autorisation, des
informations confidentielles dont ils ont eu connaissance à l'occasion
d'un audit de sécurité. »}

\medskip
\textbf{Article~69}~: \textit{« Est puni d'un emprisonnement de cinq (05) à
dix (10) ans et d'une amende de 10.000.000 (dix millions) à 100.000.000
(cent millions) F CFA ou de l'une de ces deux peines seulement, celui qui
accède sans droit, et en violation des mesures de sécurité, à l'ensemble ou
à une partie d'un réseau de communications électroniques, d'un système
d'information ou d'un équipement terminal, afin d'obtenir des informations
ou des données, en relation avec un système d'information connecté à un
autre système d'information. »}

\medskip
\textbf{Résumés} (texte officiel lu, non cité)~: article~26, alinéas~2
et~3 (mécanismes de sécurité visés par l'Agence)~; article~29, alinéa~2
(surveillance des accès aux données)~; article~30, alinéa~1 (évaluation et
révision des systèmes de sécurité)~; article~61, alinéa~3 (obstacle aux
audits)~; article~70 (attaque par saturation).
\end{boxjuris}

\section*{Règlement COBAC R-2016/04}

\begin{boxjuris}
Texte relevé sur une base réglementaire en ligne, à collationner avec la
version du Secrétariat général de la COBAC~:

\medskip
\textbf{Article~1}~: \textit{« Le présent règlement fixe les exigences
minimales relatives au système de contrôle interne que les établissements
assujettis doivent mettre en œuvre. »}

\medskip
\textbf{Article~2}, risque opérationnel~: \textit{« risque de pertes
résultant de l'inadéquation ou de la défaillance des processus internes,
des personnes et des systèmes, ou d'évènements externes. »}
\end{boxjuris}

\section*{Normes et cadres}

\begin{tableaucours}{L{4.8cm}L{3.8cm}c}{%
    \tableauHeader{Norme ou cadre} &
    \tableauHeader{Statut dans le cours} &
    \tableauHeader{Chap.}
}
ISO~31000:2018 & Paraphrasée & \ref{chap:risque}, \ref{chap:iso27005} \\
ISO/IEC~27000:2018 & Paraphrasée~; édition vérifiée & \ref{chap:grc} \\
ISO/IEC~27001:2022 et amendement~1:2024 & Paraphrasée~; éditions vérifiées
& \ref{chap:iso27000}, \ref{chap:smsi} \\
ISO/IEC~27002:2022 & Paraphrasée~; édition vérifiée~; répartition des
mesures à confirmer & \ref{chap:iso27000} \\
ISO/IEC~27004, 27005, 27007, 27014, 27017, 27018 & Paraphrasées &
\ref{chap:grc} à \ref{chap:emergents} \\
ISO~19011:2018, ISO/IEC~17021-1:2015 & Paraphrasées & \ref{chap:audit},
\ref{chap:smsi} \\
ISO/IEC~42001:2023, NIST AI RMF & Présentés & \ref{chap:emergents} \\
EBIOS Risk Manager (ANSSI) & Ateliers vérifiés sur le site de l'ANSSI~;
quatre détails à confirmer dans le guide & \ref{chap:ebios} \\
Open FAIR (O-RT 3.0.1, O-RA 2.0.1) & Versions vérifiées & \ref{chap:fair} \\
NIST CSF~2.0, SP~800-30, 800-53, 800-53A & Présentés & \ref{chap:grc},
\ref{chap:referentiels} \\
CIS Controls v8, CIS Benchmarks, PCI DSS v4.0, CVSS v4.0 & Présentés &
\ref{chap:referentiels}, \ref{chap:audit-technique} \\
\end{tableaucours}
